\documentclass[12pt, a4paper]{article}

% 宏包引入
\usepackage[UTF8]{ctex}       % 中文支持
\usepackage{amsmath, amssymb} % 数学公式与符号
\usepackage{geometry}         % 页面边距设置
\usepackage{fancyhdr}         % 页眉页脚
\usepackage{tcolorbox}        % 文本框
\tcbuselibrary{breakable}     % <--- 【关键修改1】引入跨页库
\usepackage{enumitem}         % 列表环境
\usepackage{booktabs}         % 三线表支持

% 页面布局设置
\geometry{left=2.5cm, right=2.5cm, top=2.5cm, bottom=2.5cm}

% 页眉页脚设置
\pagestyle{fancy}
\fancyhf{}
\lhead{数字信号处理（DSP）}
\rhead{公开课课前复习资料}
\cfoot{\thepage}

% 优化列表间距
\setlist[enumerate]{itemsep=0pt, parsep=0pt, topsep=5pt}
\setlist[itemize]{itemsep=0pt, parsep=0pt, topsep=5pt}

% 标题信息
\title{\textbf{数字信号处理公开课——课前复习问答清单}}
\author{23届陈文斗}
\date{2025年12月15日}

\begin{document}

	\maketitle
	\thispagestyle{fancy}

	\section*{1. 循环卷积的计算：公式法与矩阵法}

	\textbf{问题回顾：} $y_c(n) = \left[ \sum_{m=0}^{L-1} h(m)x((n-m))_L \right] R_L(n)$ 具体如何计算？

	% <--- 【关键修改2】添加 breakable 参数，允许跨页
	\begin{tcolorbox}[title=回答要点, breakable]
		循环卷积的计算在数学上和工程实现上通常有两种等效的理解视角：

		\subsection*{视角一：公式解析法（滑尺/转盘法）}
		这是公式定义的直接物理过程，适合计算特定时刻的值 $y_c(n)$。计算步骤如下：
		\begin{enumerate}
			\item \textbf{固定 (Fix)}：将序列 $h(m)$ 固定不动。
			\item \textbf{反转 (Reverse)}：将序列 $x(m)$ 以 $m=0$ 为轴进行圆周反转，得到 $x((-m))_L$。
			\item \textbf{移位 (Shift)}：根据当前要求的时刻 $n$，将反转后的 $x$ 序列右移 $n$ 位。
			\item \textbf{相乘求和 (Sum of Products)}：将固定不动的 $h(m)$ 与移位后的 $x$ 序列对齐，对应项相乘，最后将所有乘积相加得到 $y_c(n)$。
		\end{enumerate}

		\subsection*{视角二：矩阵乘法（代数法）}
		这是课堂上常用的整体计算方法，适合推导性质或编程实现。循环卷积可以表示为 $\mathbf{y} = \mathbf{X}\mathbf{h}$ 的形式：
		\[
		\begin{bmatrix} y_c(0) \\ y_c(1) \\ \vdots \\ y_c(L-1) \end{bmatrix} =
		\underbrace{
			\begin{bmatrix}
				x(0) & x(L-1) & \dots & x(1) \\
				x(1) & x(0) & \dots & x(2) \\
				\vdots & \vdots & \ddots & \vdots \\
				x(L-1) & x(L-2) & \dots & x(0)
			\end{bmatrix}
		}_{\text{循环矩阵 (Circulant Matrix)}}
		\begin{bmatrix} h(0) \\ h(1) \\ \vdots \\ h(L-1) \end{bmatrix}
		\]
		\textbf{注：} 矩阵的第一行即为 $x(n)$ 的循环倒相序列；随后的每一行均由上一行循环右移一位得到。矩阵行的每一个元素与向量 $\mathbf{h}$ 对应元素的乘积之和，即对应了视角一中的“相乘求和”。
	\end{tcolorbox}

	\section*{2. $x(n)$ 的循环倒相序列特性}

	\textbf{问题回顾：} $x(n)$ 的循环倒相序列是一个什么样的序列？

	\begin{tcolorbox}[title=回答要点, breakable]
		它是将序列 $x(n)$ 在 $0 \sim L-1$ 的区间内，关于 $n=0$ 点进行\textbf{圆周翻转}得到的序列。
		\begin{itemize}
			\item \textbf{数学表达：} $y(n) = x((-n))_L$。具体为：
			\[
			y(n) = \begin{cases}
				x(0), & n=0 \\
				x(L-n), & 1 \le n \le L-1
			\end{cases}
			\]
			\item \textbf{直观理解：} 首位 $x(0)$ 保持不变，其余点 $x(1)$ 到 $x(L-1)$ 按\textbf{逆序}排列。这也是构建循环矩阵第一行的基础。
		\end{itemize}
	\end{tcolorbox}

	\section*{3. 循环卷积矩阵的构成规律}

	\textbf{问题回顾：} 循环卷积矩阵 $x((n-m))_L$ 各行是怎么组成的？

	\begin{tcolorbox}[title=回答要点, breakable]
		该矩阵是一个\textbf{循环矩阵}（Circulant Matrix）。
		\begin{itemize}
			\item \textbf{核心规律：} 矩阵的每一行都是由\textbf{上一行向右循环移位一位}生成的。
			\item \textbf{结构分解：}
			\begin{itemize}
				\item \textbf{第 1 行：} 循环倒相序列 $[x(0), x(L-1), x(L-2), \dots, x(1)]$。
				\item \textbf{第 2 行：} 将第 1 行循环右移一位，即 $[x(1), x(0), x(L-1), \dots, x(2)]$。
				\item \textbf{第 k 行：} 对应于 $y_c(k-1)$ 的计算所需的移位序列。
			\end{itemize}
		\end{itemize}
	\end{tcolorbox}

	\section*{4. 频域采样定理恢复条件}

	\textbf{问题回顾：} 如果序列 $x(n)$ 的长度为 $M$，当频域采样点数 $N$ 满足什么条件时，时域才能恢复出 $x(n)$？

	\begin{tcolorbox}[title=回答要点, breakable]
		条件是：
		\[ \mathbf{N \ge M} \]
		\textbf{解析：}
		这是频域采样定理的核心条件。
		\begin{itemize}
			\item 频域采样会导致时域信号产生以 $N$ 为周期的周期性延拓。
			\item 只有当采样点数 $N$ 大于等于原序列长度 $M$ 时，各周期延拓分量之间才不会发生重叠。
			\item 只有不发生\textbf{时域混叠}（Time-domain Aliasing），才能通过 IDFT 完整无失真地恢复原序列 $x(n)$。
		\end{itemize}
	\end{tcolorbox}

	\section*{5. 利用 DFT 计算线性卷积的条件}

	\textbf{问题回顾：} 如果序列 $x(n)$ 长为 $M$，$h(n)$ 长为 $N$，循环卷积长为 $L$，什么条件下循环卷积等于线性卷积？

	\begin{tcolorbox}[title=回答要点, breakable]
		条件是：
		\[ \mathbf{L \ge M + N - 1} \]
		\textbf{解析：}
		\begin{itemize}
			\item 长度为 $M$ 和 $N$ 的两个序列，其\textbf{线性卷积}的非零区间长度为 $M + N - 1$。
			\item \textbf{循环卷积}实际上是线性卷积结果以 $L$ 为周期的混叠结果。
			\item 为了让循环卷积的一个周期内包含完整的线性卷积结果，必须保证周期 $L$ 足够长，使得混叠不发生。
			\item \textbf{工程实现：} 通常通过对 $x(n)$ 和 $h(n)$ 进行\textbf{补零}（Zero-padding）至长度 $L$ 来满足此条件，这也是 FFT 加速卷积运算的基础。
		\end{itemize}
	\end{tcolorbox}

\end{document}
